\section{Zadanie 4}
Niech zmienna losowa \(X_{i}^{(m)}\) oznacza liczbę kul w \(i\)-tym kubełku, przy założeniu, że rzucamy
\(m\) kul do \(n\) kubełków w sposób jednostajny i niezależny. Niech zmienne \(Y_{1}^{(m)},\dots,Y_{n}^{(m)}\)
będą niezależnymi zmiennymi z rozkładu Poissona o średniej \(\frac{m}{n}\). Podczas wykładu widzieliśmy,
że dla każdej nieujemnej funkcji \(f\) zachodzi
\[E\left(f\left(Y_{1}^{(m)},\dots,Y_{n}^{(m)}\right)\right) \geqslant E\left(f\left(X_{1}^{(m)},\dots,X_{n}^{(m)}\right)\right)P\left(\sum_{i=1}^{n}Y_{i}^{(m)}=m\right).\]
\begin{enumerate}[label=(\roman*)]
	\item Pokaż, że jeśli \(E(f(X_{1}^{(m)},\dots,X_{n}^{(m)}))\) jest rosnąca względem \(m\), to przy założeniu nieujemności \(f\) zachodzi:
	      \[E\left(f\left(Y_{1}^{(m)},\dots,Y_{n}^{(m)}\right)\right) \geqslant E\left(f\left(X_{1}^{(m)},\dots,X_{n}^{(m)}\right)\right)P\left(\sum_{i=1}^{n}Y_{i}^{(m)}\geqslant m\right).\]
	\item Wykaż, że jeśli \(E(f(X_{1}^{(m)},\dots,X_{n}^{(m)}))\) jest malejąca względem \(m\), t
	      \[E\left(f\left(X_{1}^{(m)},\dots,X_{n}^{(m)}\right)\right) \leqslant 2\cdot E\left(f\left(Y_{1}^{(m)},\dots,Y_{n}^{(m)}\right)\right).\]
\end{enumerate}
\subsection*{Rozwiązanie}
\begin{enumerate}[label=(\roman*)]
	\item
	      \[E\left(f\left(Y_{1}^{(m)},\dots,Y_{n}^{(m)}\right)\right) = \sum_{k=0}E\left(f\left(Y_{1}^{(m)},\dots,Y_{n}^{(m)}\right)|\sum_{i=1}^{n}Y_{i}=k\right)P\left(\sum_{i=1}^{n}Y_{i}^{(m)}=k\right) =\]
	      \[= \sum_{k=0}E\left(f\left(X_{1}^{(k)},\dots,X_{n}^{(k)}\right)\right)P\left(\sum_{i=1}^{n}Y_{i}^{(m)}=k\right) \geqslant \sum_{k=m}E\left(f\left(X_{1}^{(k)},\dots,X_{n}^{(k)}\right)\right)P\left(\sum_{i=1}^{n}Y_{i}^{(m)}=k\right)\]
	      Z tego że funkcja jest rosnąca mamy że:
	      \[\sum_{k=m}E\left(f\left(X_{1}^{(k)},\dots,X_{n}^{(k)}\right)\right)P\left(\sum_{i=1}^{n}Y_{i}^{(m)}=k\right) \geqslant \sum_{k=m}E\left(f\left(X_{1}^{(m)},\dots,X_{n}^{(m)}\right)\right)P\left(\sum_{i=1}^{n}Y_{i}^{(m)}=k\right) = \]
	      \[ = E\left(f\left(X_{1}^{(m)},\dots,X_{n}^{(m)}\right)\right)P\left(\sum_{i=1}^{n}Y_{i}^{(m)}\geqslant m\right)\]
	\item Tak samo jak w poprzednim możemy udowodnić że
	      \[E\left(f\left(Y_{1}^{(m)},\dots,Y_{n}^{(m)}\right)\right) \geqslant E\left(f\left(X_{1}^{(m)},\dots,X_{n}^{(m)}\right)\right)P\left(\sum_{i=1}^{n}Y_{i}^{(m)}\leqslant m\right)\]
	      \[E\left(f\left(X_{1}^{(m)},\dots,X_{n}^{(m)}\right)\right) \leqslant \frac{E\left(f\left(Y_{1}^{(m)},\dots,Y_{n}^{(m)}\right)\right)}{P\left(\sum_{i=1}^{n}Y_{i}^{(m)}\leqslant m\right)}\]
	      \[P\left(\sum_{i=1}^{n}Y_{i}^{(m)}\leqslant m\right) \geqslant \frac{1}{2}\].
	      \[E\left(f\left(X_{1}^{(m)},\dots,X_{n}^{(m)}\right)\right) \leqslant 2E\left(f\left(Y_{1}^{(m)},\dots,Y_{n}^{(m)}\right)\right)\]
\end{enumerate}